% Spectral structure used in the intrinsic five-generator bound.
\subsubsection{Maximal idempotents and the weight operator}
\label{subsec:nr-maximal-idempotents}

By \cref{lem:square-free-critical-set}, an irreducible harmonic $F$
has a nonsingular real zero. Thus \cref{prop:divisibility}
gives $\cL(F)=-qF$, where homogeneity makes $q$ quadratic.

Use the commutative Euclidean metrized algebra and normalization of
\cref{prop:cubic-multiplication-dictionary}.
Harmonicity is the condition $\tr L_x=0$ for every $x$;
we call such an algebra \emph{exact}, following
\cite[Section~3.2 and Definition~4.1]{FoxTkachev24}.
For a self-adjoint operator \(T\), the equation
\(\cL(F)=-\frac32\ip{Tx}{x}F\) is then equivalent to
\begin{equation}\label{eq:nr-Hsiang-identity}
 \ip{x^2}{x^3}=\ip{Tx}{x}\ip{x^2}{x}.
\end{equation}
This defines an \emph{exact Hsiang algebra}, with symmetric bilinear weight
$b(x,y)=\ip{Tx}{y}$. The radial case is the special case in which
$T$ is scalar.
Here \(x^3=x(x^2)\) and \(x^4=x(x^3)\); no associativity is assumed.
Differentiating \eqref{eq:nr-Hsiang-identity} gives
\begin{equation}\label{eq:nr-vector-identity}
 4x^4+(x^2)^2-2\ip{x^2}{x}Tx-3\ip{Tx}{x}x^2=0.
\end{equation}

As shown in \cref{subsec:cubic-multiplication},
$\operatorname{Ann}(A)=\{v:vA=0\}$ is the cylindrical subspace of $F$.
Metric invariance makes $A_{\mathrm{ess}}=\operatorname{Ann}(A)^\perp$
a subalgebra; its orthogonal splitting from the annihilator has zero
mixed products.

We use the spherical maximum construction of an idempotent
\cite[Lemma~6.6.1]{NadirashviliTkachevVladut14} and derive the spectral
bounds below.

\begin{lemma}[Maximal normalization]\label{lem:nr-max-idempotent}
Let \(A\) be a commutative metrized algebra with a Euclidean metric
and nonzero multiplication. After multiplying its product
by a positive constant, there is a unit idempotent \(c\) such that
\begin{equation}\label{eq:nr-spherical-maximum}
 \max_{|x|=1}\ip{x^2}{x}=\ip{c^2}{c}=1.
\end{equation}
Every eigenvalue of \(L_c\) on \(c^\perp\) lies in
\([-1,\frac12]\).
\end{lemma}

\begin{proof}
Put \(h(x)=\ip{x^2}{x}\). Polarization shows that \(h\) is a
nonzero odd cubic, so its maximum \(M\) on the unit sphere is positive.
At a maximizer \(c\), the Lagrange multiplier equation is \(c^2=Mc\).
Replacing the product by \(xy/M\) proves
\eqref{eq:nr-spherical-maximum}. It preserves exactness and replaces
\(T\) in \eqref{eq:nr-Hsiang-identity} by \(T/M^2\).
The constrained Hessian of \(h\) at \(c\) is
\(6L_c-3\Id\) on \(c^\perp\), giving the upper bound.
For a unit eigenvector \(u\perp c\), with \(cu=ru\), maximality gives
\[
 h\left(-\frac c2\pm\frac{\sqrt3}2u\right)
 =-\frac18-\frac98r\pm\frac{3\sqrt3}8h(u)\le1.
\]
Averaging gives \(r\ge-1\).
\end{proof}

Fix this normalization, which rescales the cubic by a positive constant.

\begin{lemma}[The weight operator at a unit idempotent]
\label{lem:nr-idempotent-quotient}
Every unit idempotent \(e\) in an algebra satisfying
\eqref{eq:nr-Hsiang-identity} satisfies
\begin{equation}\label{eq:nr-quotient-polynomial}
 Te=e,\qquad T=p(L_e)-6P_e,\qquad p(t)=4t^3+4t^2-t,
\end{equation}
where \(P_ex=\ip{e}{x}e\). In particular, \([T,L_e]=0\).
\end{lemma}

\begin{proof}
At \(e\), the scalar identity gives \(\ip{Te}{e}=1\), and the
vector identity gives \(Te=e\). Its derivative at \(e\) is
\[
 8L_e^3y+8L_e^2y-2L_ey-2Ty-12\ip{e}{y}e=0.
\]
This proves \eqref{eq:nr-quotient-polynomial}. Since \(L_e\) is
self-adjoint and \(L_ee=e\), it commutes with \(P_e\).
\end{proof}

Write \(E_r=\ker(L_c-r\Id)\). By maximality \(E_1=\R c\), and
\begin{equation}\label{eq:nr-Peirce-quotient}
 Tu=p(r)u\qquad (u\in E_r\subset c^\perp).
\end{equation}

\subsubsection{Fusion identities and spectral rigidity}
\label{subsec:nr-spectral-rigidity}

We adapt the radial polarization argument in
\cite[Lemma~6.6.2]{NadirashviliTkachevVladut14}, specifically equation~(6.6.8).
For nonradial $T$, we derive the fusion rule directly from
\eqref{eq:nr-vector-identity}, using
\begin{equation}\label{eq:nr-fusion-polynomial}
 \mathscr K(r,s,t)=4(r^2+s^2+t^2+rs+rt+st)+2(r+s+t)-3.
\end{equation}

\begin{lemma}[Fusion and the square identity]\label{lem:nr-fusion-square}
For \(u\in E_r\), \(v\in E_s\), \(w\in E_t\), all
perpendicular to \(c\), one has
\begin{equation}\label{eq:nr-fusion}
 \mathscr K(r,s,t)\ip{uv}{w}=0.
\end{equation}
For \(u\in E_r\subset c^\perp\), set
\(Q_r(t)=\mathscr K(r,r,t)\). Then
\begin{equation}\label{eq:nr-square-operator}
 Q_r(L_c)u^2=3\bigl(2r+p(r)\bigr)|u|^2c.
\end{equation}
Writing \(u^2=r|u|^2c+\sum_tq_t\), with
\(q_t\in E_t\cap c^\perp\), gives
\begin{equation}\label{eq:nr-square-norm}
 \sum_t(t+4r)|q_t|^2=4r^2(3r-1)(r+1)|u|^4.
\end{equation}
\end{lemma}

\begin{proof}
The mixed coefficient of \(ab\) in
\eqref{eq:nr-vector-identity} at \(c+au+bv\), projected to
\(c^\perp\) and divided by two, is the projection of
\[
 \bigl[4L_c^2+(4r+4s+2)L_c
       +4(r^2+rs+s^2)+2(r+s)-3\bigr]uv.
\]
Pairing with \(w\) proves \eqref{eq:nr-fusion}. The coefficient
of \(a^2\) at \(c+au\) gives
\eqref{eq:nr-square-operator}; thus \(Q_r(t)q_t=0\).
For the norm identity, pair the vector identity at \(u\) with \(c\).
Metric invariance gives
\[
 \ip{u^4}{c}=r|u^2|^2,\qquad
 \ip{(u^2)^2}{c}=r^2|u|^4+\sum_t t|q_t|^2.
\]
The term containing \(Tu\) pairs to zero, and the last term is
\(-3p(r)r|u|^4\). Rearranging yields
\eqref{eq:nr-square-norm}, since
\(3p(r)r-4r^3-r^2=4r^2(3r-1)(r+1)\).
\end{proof}

\begin{proposition}[A uniform spectral interval]
\label{prop:nr-spectral-interval}
Put \(r_0=-(2+\sqrt7)/6\). Then
\begin{equation}\label{eq:nr-Peirce-interval}
 \operatorname{Spec}(L_c|_{c^\perp})
 \subset\{-1,0\}\cup[r_0,-1/6]\cup[1/3,1/2].
\end{equation}
Moreover,
\begin{equation}\label{eq:nr-quotient-interval}
 \operatorname{Spec}T\subset\{0\}\cup
 \left[\frac7{27},\frac{17+7\sqrt7}{27}\right],
 \qquad \ker T=\operatorname{Ann}(A).
\end{equation}
\end{proposition}

\begin{proof}
All possible output eigenvalues \(t\) in
\eqref{eq:nr-square-norm} lie in \([-1,1/2]\). If
\(-1<r<r_0\), then \(Q_r(t)=4t^2+(8r+2)t+12r^2+4r-3\)
is decreasing on that interval and
\(Q_r(1/2)=12r^2+8r-1>0\). All \(q_t\) therefore vanish,
contradicting the negative right side of
\eqref{eq:nr-square-norm}. If \(-1/6<r<0\), both endpoint
values
\[
 Q_r(-1)=(6r+1)(2r-1),\qquad Q_r(1/2)=12r^2+8r-1
\]
are negative. Convexity makes \(Q_r\) negative throughout the
interval, giving the same contradiction. If \(0<r<1/3\),
then \(Q_r(-1)<0\) and \(Q_r(0)<0\); hence every admissible
root is positive. The left side of the norm identity is nonnegative
and the right side is negative. This proves
\eqref{eq:nr-Peirce-interval}. At \(r=1/3\), the norm identity
also shows that every \(q_t\) vanishes, because \(t+4/3>0\).

Applying \(p\) to these intervals gives the spectral inclusion in
\eqref{eq:nr-quotient-interval}. The only zero of \(p\) allowed
by \eqref{eq:nr-Peirce-interval} is \(0\). Thus \(u\in\ker T\)
implies \(cu=0\), and \eqref{eq:nr-square-operator} gives
\[
 (4L_c^2+2L_c-3\Id)u^2=0.
\]
Its two roots \((-1\pm\sqrt{13})/4\) lie outside
\([-1,1/2]\), and it is nonzero at \(1\). Consequently \(u^2=0\).
The coefficient of \(t^3\) in the vector identity at \(x+tu\)
is then \(8L_u^3x\): all other terms have degree at most two,
since \(Tu=u^2=0\). Hence \(L_u^3=0\), and self-adjointness
gives \(L_u=0\). Conversely, an annihilator element is orthogonal
to \(c=c^2\) and is killed by \(L_c\), so
\eqref{eq:nr-quotient-polynomial} kills it as well.
\end{proof}

By \cref{prop:nr-spectral-interval}, restriction to $A_{\mathrm{ess}}$
preserves exactness and the defining identity, and
$T|_{A_{\mathrm{ess}}}$ is positive definite.

Let
\begin{equation}\label{eq:nr-eigenvalue-one-space}
 B=E_T(1)=\R c\oplus Z\oplus H\oplus W,
 \qquad Z=E_{-1},\quad H=E_{-1/2},\quad W=E_{1/2},
\end{equation}
since \(p(t)-1=(t+1)(2t-1)(2t+1)\).
The fusion polynomial gives the following possible outputs on
\(c^\perp\):
\[
\begin{array}{c|cccccc}
(r,s)&(-1,-1)&(-1,-\frac12)&(-1,\frac12)&
(-\frac12,-\frac12)&(-\frac12,\frac12)&(\frac12,\frac12)\\ \hline
 t&\text{none}&\frac12&-\frac12,\frac12&-\frac12&
 -1,\frac12&-1,-\frac12
\end{array}
\]
(The additional root \(1\) in the fourth column is absent on
\(c^\perp\).) In particular,
\begin{equation}\label{eq:nr-B-subalgebra}
 B^2\subset B,\qquad BB^\perp\subset B^\perp.
\end{equation}

A Euclidean commutative metrized algebra is \emph{eikonal} if
\(x^3=|x|^2x\), equivalently \(|x^2|^2=|x|^4\).
Equivalently, the associated cubic $6F$ is a cubic eiconal in the sense
of \cite[Definition~3.1.1]{NadirashviliTkachevVladut14}; exactness is a
separate condition.

\begin{lemma}[The negative Peirce spaces]\label{lem:nr-negative-spaces}
\begin{enumerate}[label=\textup{(\roman*)}]
\item If \(H\ne0\), then \(\operatorname{Spec}T\subset\{1,0\}\).
\item If \(Z\ne0\), then
\(\operatorname{Spec}T\subset\{1,7/27,0\}\).
For \(U=E_T(7/27)\), one has
\begin{equation}\label{eq:nr-U-peirce}
 U=P\oplus Q,\quad P=E_{1/3},\quad Q=E_{-1/6},\qquad
 L_z|_P=0,\quad (L_z|_Q)^2=\frac1{12}|z|^2\Id_Q.
\end{equation}
\end{enumerate}
\end{lemma}

The proof is given in \cref{app:nr-spectral-exclusions}.

Exactness excludes $Z=H=W=0$.

\begin{lemma}\label{lem:nr-not-one-dimensional}
In an exact Hsiang algebra whose multiplication is not identically zero,
\(E_T(1)\ne\R c\).
\end{lemma}

The proof is given in \cref{app:nr-spectral-exclusions}.

\begin{theorem}[Spectral rigidity]\label{thm:nonradial-spectral}
Every exact Euclidean Hsiang algebra whose multiplication is not
identically zero admits a positive rescaling for which
\begin{equation}\label{eq:nr-rigid-spectrum}
 \operatorname{Spec}T\subset\{0,7/27,1\},\qquad
 \ker T=\operatorname{Ann}(A).
\end{equation}
On $A_{\mathrm{ess}}$, either $T|_{A_{\mathrm{ess}}}=\Id$ or
\begin{equation}\label{eq:nr-quotient-splitting}
 A_{\mathrm{ess}}=B\oplus U,\qquad T|_B=\Id,\qquad T|_U=\frac7{27}\Id.
\end{equation}
\end{theorem}

\begin{proof}
The kernel assertion is \cref{prop:nr-spectral-interval}.
If \(H\ne0\) or \(Z\ne0\), use
\cref{lem:nr-negative-spaces}.
Otherwise the fusion table gives \(w^2=|w|^2c/2\) for \(w\in W\).
The norm identity with \(r=1/2\) then reads
\(0=3|w|^4/4\), so \(W=0\). This would give
\(B=\R c\), contrary to \cref{lem:nr-not-one-dimensional}.
\end{proof}
